\chapter{Discussion}
\label{ch:discussion}

\section{One account of the results}
\label{sec:story}

The results fit a single account, and each part of it was tested rather than
assumed.

\begin{enumerate}
  \item Instruction tuning teaches a RAG model to defer to its context, and
    above all to decline when the context looks inadequate. This shows up
    without any attack, as a 29--37\% refusal floor on answerable questions
    (\cref{sec:res-floor}). That floor rises steadily as retrieval quality
    falls (\cref{fig:gradient}).
  \item The strongest attacks do not inject unsafe content or exhaust compute.
    They make the evidence look circular, contradictory or crowded out
    (\cref{sec:res-attacks}). A faithful model then does what it was trained to
    do, and declines.
  \item The deference is operative because of training, not wording. The same
    refusal licence moves an aligned model 3.1 times as far as its base model
    (\cref{sec:res-prompt}).
  \item This is not inherent to post-training. Reasoning distillation reduces
    the deference and reverses the effect (\cref{sec:res-r1}).
  \item The base model's apparent robustness is mostly a failure to notice.
    The same attacks cost it fewer answers, but those answers become silent
    errors (\cref{sec:res-correctness}).
\end{enumerate}

The contribution is therefore not ``alignment is bad for RAG''. It is that
alignment changes \emph{which} property an attacker on the corpus can degrade,
from integrity to availability. A deployer choosing between the two models is
choosing which failure they would rather detect.

\todo{Expand: implications for (a) how over-refusal is evaluated (content-only,
no unsafe text), (b) RAG evaluation that reports only accuracy, (c) the defence
result that reranking is not free.}

\section{Threats to validity}
\label{sec:threats}

\paragraph{Prompt confound.} Every \texttt{auto}-prompt contrast is an upper
bound. The prompt-matched estimate ($+0.134$ floor, $+0.204$ ASR) comes from one
pair, one attack and one corpus (\cref{sec:res-prompt}). Two results are immune
to the confound regardless: the R1 pair, where both sides receive the same
prompt, and the within-model retrieval gradient, which holds the prompt fixed.

\paragraph{The withdrawn corpus.} HotpotQA was run on 12 cells and then
withdrawn. Its corpus builder kept one gold passage per query, but every
HotpotQA question needs exactly two. Half the evidence for every query was
therefore never indexed, and the aligned model's 73--81\% refusal rate there is
largely \emph{correct}. Its numbers, which were the largest effects in the
project, are excluded throughout. FiQA, screened for this defect before use,
replaces it. This was the second time an audit of the benchmark's own
instruments changed a headline. The first was the refusal classifier. Three defects
found in it were fixed before the validation in \cref{sec:classifier}, and
fixing them moved the headline from 12 genuine paradoxes in 50 runs to 39 in
62.

\paragraph{Selection in ASR.} Attributable ASR conditions on queries answered
clean. As the floor rises, the survivors are increasingly the queries the model
was willing to answer despite poor retrieval, a self-selected robust subset.
Denominators are reported, and the unconditional measure, which has the
opposite bias, is reported alongside it.

\paragraph{Instrument.} The classifier is validated (\cref{tab:classifier}), but
its gold set is not fully independent of its author. Residual error is
asymmetric towards missed refusals, which biases the gap \emph{down}.

\paragraph{Scope.} Models are 7--8B parameters, 4-bit quantised, with a single
small embedder and $k=5$. Correctness covers 656 NQ queries only. Variance was
measured on only one condition (C1 at 5\%, 1.5 points). Generation is greedy, so
the main grid is deterministic.

\todo{Larger models; other embedders; an adaptive attacker.}
